% ==============================================================================
% LISTA DE ABREVIATURAS E SIGLAS (opcional; NBR 14724:2024)
% - Ordem ALFABÉTICA.
% - Incluir apenas as siglas efetivamente usadas no texto. Na primeira
%   ocorrência, a sigla vem entre parênteses, após o nome por extenso:
%   "Comitê de Ética em Pesquisa (CEP)". Depois, usa-se só a sigla.
% - Siglas estrangeiras: por extenso no idioma original, em itálico, e,
%   se útil, a tradução entre parênteses.
% - Não entra no sumário.
% ==============================================================================
\begin{listaabreviaturas}
  CEP    & Comitê de Ética em Pesquisa \\
  CNS    & Conselho Nacional de Saúde \\
  CONEP  & Comissão Nacional de Ética em Pesquisa \\
  IA     & Inteligência Artificial \\
  IAGen  & Inteligência Artificial Generativa \\
  OSF    & \textit{Open Science Framework} \\
  POC    & Prova de conceito (\textit{proof of concept}) \\
  PRISMA & \textit{Preferred Reporting Items for Systematic Reviews and Meta-Analyses} \\
  SHAP   & \textit{SHapley Additive exPlanations} \\
  XAI    & \textit{Explainable Artificial Intelligence} (Inteligência Artificial Explicável) \\
\end{listaabreviaturas}

% ==============================================================================
% LISTA DE SÍMBOLOS (opcional; NBR 14724:2024)
% - Ordem em que os símbolos APARECEM NO TEXTO (e não alfabética).
% - Cada símbolo com o seu significado; um símbolo, um único significado em
%   todo o trabalho (por isso a partição de validação usa o índice m, e não
%   j, já reservado ao jogador/variável na Equação 2.1).
% - Atualizar sempre que uma equação for incluída, alterada ou removida.
% - Não entra no sumário.
% - Cada linha: \simbolo{símbolo}{significado}. O comando não usa tabela,
%   o que evita conflitos de pacotes com colunas de largura fixa.
% ==============================================================================
\providecommand{\simbolo}[2]{%
  \noindent\parbox[t]{2.2cm}{#1}\parbox[t]{13.2cm}{#2}\par\smallskip}

\begin{listasimbolos}
  \simbolo{$j$}{Índice de um jogador (variável preditora)}
  \simbolo{$\phi_j(v)$}{Valor de Shapley do jogador $j$}
  \simbolo{$v$}{Função característica, que associa a cada coalizão o seu ganho}
  \simbolo{$S$}{Coalizão de jogadores que não contém $j$}
  \simbolo{$N$}{Conjunto de jogadores (variáveis preditoras)}
  \simbolo{$|N|$}{Cardinalidade do conjunto $N$ (número de elementos)}
  \simbolo{$VP$}{Número de verdadeiros positivos}
  \simbolo{$FP$}{Número de falsos positivos}
  \simbolo{$FN$}{Número de falsos negativos}
  \simbolo{$\tau$}{Limiar de decisão do classificador}
  \simbolo{$\hat{p}_i$}{Probabilidade estimada da classe positiva para a observação $i$}
  \simbolo{$i$}{Índice de uma observação}
  \simbolo{$F_1$}{Média harmônica entre precisão e revocação (\textit{F1-score})}
  \simbolo{$\theta$}{Vetor de hiperparâmetros de um modelo}
  \simbolo{$k$}{Número de partições da validação cruzada}
  \simbolo{$\hat{\theta}$}{Hiperparâmetros selecionados por validação cruzada}
  \simbolo{$\argmin$}{Argumento que minimiza a expressão}
  \simbolo{$\Theta$}{Espaço de busca dos hiperparâmetros}
  \simbolo{$\mathcal{L}$}{Função de perda}
  \simbolo{$D_m$}{$m$-ésima partição de validação}
  \simbolo{$m$}{Índice de uma partição de validação}
\end{listasimbolos}
